\documentclass[pdflatex,sn-mathphys-num]{sn-jnl} % Math and Physical Sciences Numbered Reference Style
\usepackage[T1]{fontenc}
\usepackage{lmodern}
\input{glyphtounicode}
\pdfgentounicode=1
\usepackage[numbers,sort&compress]{natbib}
\usepackage{amsmath,amssymb,amsfonts}
\usepackage{graphicx}
\usepackage{hyperref}
\usepackage{physics}
\usepackage{tikz}
\usepackage{amsopn}
\usepackage{amstext}
\usepackage{microtype}
\usepackage{xparse}
\usepackage{algorithm}
\usepackage{algorithmicx}
\usepackage{doi}
\usepackage{algpseudocode}
% Boxes
\usepackage[most]{tcolorbox}
% Plots (pgfplots provides axis/addplot)
\usepackage{pgfplots}
\usetikzlibrary{intersections}
\usepgfplotslibrary{fillbetween}
\pgfplotsset{compat=1.18}
\NewDocumentCommand{\subsubsubsection}{s m}{%
  \IfBooleanTF{#1}{\paragraph*{#2}}{\paragraph{#2}}%
}
\usetikzlibrary
{arrows.meta, calc, fit, decorations.pathmorphing, decorations.pathreplacing, shapes.geometric, shadings, positioning,
  decorations.markings, patterns}
\theoremstyle{remark}
\newtheorem*{remark}{Remark}
\theoremstyle{plain}
\newtheorem{theorem}{Theorem}
\newtheorem{lemma}{Lemma}
\newtheorem{corollary}[theorem]{Corollary}
\newtheorem{hypothesis}{Hypothesis}

\title{Conditions for an Infrared Einstein Sector from Spectral Geometry}
\author*[1]{\fnm{Jérôme} \sur{Beau}}\email{jerome.beau@cosmochrony.org}
\affil*[1]{\orgname{Independent Researcher}, \country{France}}
\date{}
\hypersetup{
  pdftitle={Conditions for an Infrared Einstein Sector from Spectral Geometry},
  pdfauthor={Jérôme Beau},
  pdfsubject={Conditions for a local infrared Einstein sector from a four-dimensional Laplace-type operator},
  pdfkeywords={spectral geometry, heat-kernel expansion, spectral cutoff, zeta regularization, effective action,
    Einstein tensor, Born--Infeld completion, renormalization matching}
}

\abstract{
  We determine the conditions under which a four-dimensional Laplace-type operator on a supplied Riemannian metric
  can support a local infrared Einstein sector.
  The analysis separates two regularization statements that must not be conflated.
  A physical proper-time cutoff $\Lambda=\ell_{\mathrm{sp}}^{-1}$ produces power-sensitive local terms
  $a_0\Lambda^4$ and $a_2\Lambda^2$, while zeta regularization produces finite and logarithmic scale dependence
  governed in four dimensions by $a_4=\zeta_A(0)$.

  For $S_\Pi=+\tfrac12\log\det' A_g$, one minimal scalar contributes
  $\Delta c_{\mathrm{EH}}=-\Lambda^2/[12(4\pi)^2]$ to the Einstein--Hilbert coefficient.
  The finite renormalized value, its sign, and the observed Newton constant require a matching condition and the
  complete operator content; they are not predicted by this determinant alone.
  Once these coefficients are supplied, the metric variation decomposes into Einstein, cosmological,
  higher-derivative, and non-local sectors.
  The Einstein term dominates locally when the matched Einstein coefficient is non-zero and the renormalized
  coefficient ratio satisfies $R L_4^2\ll1$, a power-counting statement that does not bound the non-local terms.

  Under the separate coherence-extensivity hypothesis, the determinantal Born--Infeld density is an admissible
  tensorial completion of a supplied Einstein infrared term, but the structural conditions do not select it uniquely:
  an explicit one-parameter family of distinct admissible completions shares its infrared limit and saturation
  behaviour; among linear two-derivative carriers, the saturation condition on the Ricci endomorphism selects that
  endomorphism as the spectral carrier.
  \textbf{Interpretive status.}
  Spectral geometry fixes the allowed local structures and their cutoff sensitivities, but it does not yet derive the
  positive finite gravitational coupling.
  The result identifies the matching problem that must be solved before gravity can be claimed as an emergent spectral
  prediction.
}

\keywords
{Spectral geometry, Heat-kernel expansion, Spectral cutoff,
  Zeta regularization, Effective action, Einstein tensor,
  Born--Infeld completion, Renormalization matching}

\begin{document}
  \maketitle

  \tableofcontents

  \clearpage

  \section{Introduction}
    \label{sec:introduction}

    Gravity is traditionally described as a geometric phenomenon.
    In general relativity, the gravitational interaction is encoded in the curvature of spacetime, and the Newtonian
    potential arises as a weak-field limit of the Einstein equations.
    This geometric formulation is remarkably successful.
    It nevertheless leaves open the conceptual question of whether gravitational dynamics must be fundamental, or
    whether it may arise as an effective response of a more primitive operator-based structure~\cite{Cosmochrony}.

    In parallel, operator-theoretic and spectral approaches have clarified how elliptic operators encode geometric and
    physical information.
    Functional determinants, resolvents, and heat-kernel asymptotics play a central role in quantum field theory,
    statistical mechanics, and spectral geometry~\cite{Vassilevich2003,ParkerToms2009}.
    Variations of logarithmic determinants generate non-local responses through the inverse operator.
    This suggests a bridge between entropy-like functionals and long-range interaction.

    In a previous analysis~\cite{Beau2026a}, we investigated a minimal mechanism by which a Newtonian interaction
    can emerge from the
    variation of a projective entropy functional in three spatial dimensions.
    For a positive Laplace-type operator \(A\), the functional
    \begin{equation}
      S_\Pi[A] = \frac12 \log \det\nolimits' A
    \end{equation}
    has the weak-perturbation variation
    \begin{equation}
      \delta S_\Pi = \frac12 \Tr\!\left(A^{-1}\delta A\right),
    \end{equation}
    so that the Green kernel governs the spatial structure of the response.
    In three dimensions, the Green kernel of a Laplace-type operator exhibits a universal \(1/r\) decay.
    Consequently, the entropy response inherits a Newtonian long-range profile without assuming a force law.

    The goal of the present paper is to develop a covariant extension of this mechanism and to identify, in a controlled
    infrared regime, the leading local geometric content of the \emph{renormalized} metric variation.
    We do not claim an exact equality between the metric variation of the bare determinant and Einstein's equations.
    We also do not discard non-local contributions or higher-derivative logarithmic terms.
    Instead, our central claim is the following.

    \paragraph{Central claim (scheme separation and conditional IR locality).}
      In four dimensions, the $a_2$ coefficient produces a quadratic Einstein--Hilbert contribution only in an
      explicit cutoff scheme, while zeta-scale dependence is governed by $a_4$.
      The finite renormalized Einstein coefficient is a matching datum.
      If that coefficient is nonzero and the renormalized higher-derivative ratios satisfy $R L_4^2\ll1$, the
      Einstein tensor is the dominant local two-derivative metric response.
      The structural necessity of non-injectivity for any genuinely emergent description is established
      in~\cite{Beau2026l}.

      The starting point of the present paper is the Laplace-type operator $A_g = -\nabla^2_g$ of a supplied
      Riemannian metric $g$.
      Its identification with an operator constructed from the Cosmochrony admissibility data is not established.
      Q5a~\cite{Beau2026q5a} shows that the published admissibility form converges to the zero form, so that a
      non-zero spatial limit operator is the unestablished input [H-L] of Q5b~\cite{Beau2026q5b}.
      The lift of that operator to an effective operator on
      $\mathbb{R}_\tau \times \mathrm{Heis}_3(\mathbb{R})$ is the open lifting hypothesis [H-lift] (Q5b,
      Hypothesis~4.1), which Q9~\cite{Beau2026q9} does not discharge.
      Under [H-L] and [H-lift], Q5b gives a co-metric of rank three.
      Reading the supplied metric $g$ as an emergent metric would further require a supplied rank-four extension,
      and relating it to a Lorentzian metric needs a continuation that is not performed.
      The results below hold for any supplied Riemannian metric and do not depend on that reading.
      The present paper establishes the available infrared tensor structure and its regulator dependence.
      It does not establish that the minimal determinant alone fixes a positive observed Newton coupling.

      Operationally, this result identifies the coefficient ratios that control the derivative hierarchy.
      Their relation to $\ell_{\mathrm{sp}}$ is conditional on cutoff-dominated matching without cancellations.

      The remainder of the paper is organized as follows.
      In Section~\ref{sec:projective-entropy-covariant}
      , we define the covariant projective entropy functional and set up the
      minimal Laplace-type operator choice.
      In Section~\ref{sec:hk-sdw}
      , we separate the proper-time-cutoff and zeta roles of the Seeley--DeWitt coefficients.
      In Section~\ref{sec:metric-variation}
      , we define the renormalized spectral stress tensor and decompose it by derivative
      order into Einstein, quadratic-curvature, and non-local pieces.
      In Section~\ref{subsec:ir-hierarchy}
      , we formulate the infrared hierarchy as a covariant derivative expansion and
      state the dominance criterion \(R\,L_4^2\ll 1\) for a non-zero matched Einstein coefficient.
      In Section~\ref{sec:newton-identification}
      , we discuss the cutoff-sensitive contribution and the matching definition of the effective Newton coupling.
      We conclude in Section~\ref{sec:discussion} with limitations and outlook.

  \section{Covariant projective entropy functional}
  \label{sec:projective-entropy-covariant}

  \subsection{Ellipticity and minimal Laplace-type operator}
    \label{subsec:ellipticity}

    Let $(M,g)$ be a smooth four-dimensional Riemannian manifold.
    We consider a second-order differential operator acting on scalar functions.

    The analysis relies crucially on ellipticity.
    For a differential operator of order two, ellipticity means that its
    principal symbol is invertible away from the zero section.

    We choose the minimal Laplace-type operator
    \begin{equation}
      A_g = -\nabla_g^2 ,
    \end{equation}
    where $\nabla_g$ is the Levi-Civita connection.
    Its principal symbol is
    \begin{equation}
      \sigma_2(A_g)(x,k) = g^{\mu\nu}(x) k_\mu k_\nu .
    \end{equation}
    On a Riemannian manifold, $g^{\mu\nu}k_\mu k_\nu>0$ for all nonzero $k$.
    Therefore $A_g$ is elliptic.

    The principal symbol fully determines the local propagation structure.
    In particular, the metric enters exclusively through $\sigma_2(A_g)$.
    Lower-order terms do not affect ellipticity.

    Ellipticity ensures:

    \begin{itemize}
      \item existence of the resolvent $(A_g+\lambda)^{-1}$ for $\lambda>0$,
      \item discreteness of the spectrum on compact manifolds,
      \item local regularity of the Green kernel,
      \item validity of the short-time heat-kernel expansion,
      \item meromorphic continuation of the spectral zeta function.
    \end{itemize}

    More general scalar Laplace-type operators may include a curvature coupling,
    \begin{equation}
      A_g(\xi) = -\nabla_g^2 + \xi R .
    \end{equation}

    The additional term $\xi R$ modifies only the zeroth-order endomorphism
    part of the operator and does not alter its principal symbol.
    It therefore does not affect ellipticity.

    In this work we set $\xi=0$.
    This is the minimal operator choice in which the effective geometry is
    entirely encoded in the principal symbol, without additional curvature
    input introduced at lower order.

    The conformal value $\xi=1/6$ cancels the $a_2$ coefficient in four
    dimensions and suppresses the cutoff-sensitive Einstein--Hilbert contribution.
    The present analysis focuses on the minimal elliptic case $\xi=0$,
    for which the $a_2$ contribution is nonvanishing.

  \subsection{Zeta determinant and spectral-cutoff matching}
    \label{subsec:zeta-regularization}

    We define the projective entropy functional by
    \begin{equation}
      S_\Pi[g] \equiv \frac12 \log \det\nolimits' A_g ,
    \end{equation}
    where the prime denotes omission of zero modes, if any.

    Its zeta-regularized form at reference scale $\mu$ is
    \begin{equation}
      S_\Pi^{\zeta}[g;\mu]
      =
      -\frac12\zeta'_{A_g}(0)
      -\frac12\log(\mu^2)\,\zeta_{A_g}(0),
      \label{eq:zeta-functional}
    \end{equation}
    with
    \begin{equation}
      \zeta_{A_g}(s)
      =
      \frac{1}{\Gamma(s)}
      \int_0^\infty
      t^{\,s-1}
      \Tr\!\left(e^{-tA_g}\right)\,dt ,
    \end{equation}
    initially defined for $\Re(s)$ sufficiently large and continued meromorphically to $s=0$.

    In four dimensions, $\zeta_{A_g}(0)$ is governed locally by the integrated coefficient $a_4$.
    Consequently, zeta-scale dependence is logarithmic and belongs to the four-derivative sector.
    It does not generate an $a_2\mu^2R$ term.

    Cosmochrony also supplies a physical spectral resolution scale
    $\ell_{\mathrm{sp}}$.
    To track power-sensitive local terms, we therefore introduce the separate proper-time-cutoff functional
    \begin{equation}
      S_\Pi^{\Lambda}[g]
      =
      -\frac12
      \int_{\Lambda^{-2}}^\infty
      \frac{\dd t}{t}\,
      \Tr\!\left(e^{-tA_g}\right),
      \qquad
      \Lambda\equiv\ell_{\mathrm{sp}}^{-1}.
      \label{eq:cutoff-functional}
    \end{equation}
    The cutoff and zeta prescriptions answer different questions and will not be identified.
    The cutoff determines the power-sensitive matching contributions, while the zeta determinant organizes the finite
    and logarithmic sectors.

    The renormalized functional additionally contains finite local coefficients fixed by matching conditions.
    In particular, the finite value and sign of the Einstein--Hilbert coefficient are not determined by the heat-kernel
    expansion alone.

  \subsection{First variation and resolvent trace}
    \label{subsec:first-variation}

    For a perturbation $A_g \to A_g + \delta A_g$, the first variation of the determinant satisfies formally
    \begin{equation}
      \delta S_\Pi
      =
      \frac12 \Tr\!\left(A_g^{-1}\,\delta A_g\right) .
    \end{equation}

    This identity shows that the response of the functional is governed by the resolvent kernel.
    In three spatial dimensions, this mechanism produces a universal $1/r$
    Green profile, leading to a Newtonian response without assuming a fundamental force law.

    In the present covariant four-dimensional setting, the metric variation defines a renormalized spectral stress
    tensor,
    \begin{equation}
      \mathcal{T}^{\Pi,{\mathrm ren}}_{\mu\nu}
      \equiv
      -\frac{2}{\sqrt{g}}
      \frac{\delta S_\Pi^{\mathrm ren}}{\delta g^{\mu\nu}} .
    \end{equation}

    Its precise geometric decomposition depends on the heat-kernel structure of $A_g$.
    In particular, local two-derivative, local four-derivative, and non-local contributions appear with distinct
    regulator dependence.
    The classification of these terms and their infrared hierarchy is developed in the following sections.

  \section{Heat-kernel expansion and Seeley--DeWitt structure}
  \label{sec:hk-sdw}

  \subsection{Short-time heat-kernel expansion}
    \label{subsec:hk-expansion}

    For an elliptic Laplace-type operator $A_g$ on a smooth four-dimensional
    Riemannian manifold, the heat trace admits a short-time asymptotic expansion
    as $t \to 0^+$ of the form
    \begin{equation}
      \Tr\!\left(e^{-tA_g}\right)
      \sim
      \frac{1}{(4\pi t)^2}
      \sum_{k\ge 0}
      t^{k}
      \int_M d^4x\,\sqrt{g}\, a_{2k}(x) .
    \end{equation}

    The coefficients $a_{2k}(x)$ are local curvature invariants constructed from
    the metric and its derivatives.
    They are the Seeley--DeWitt coefficients of the operator~\cite{Seeley1967,DeWitt1965,Vassilevich2003}.

    For the minimal scalar Laplacian ($\xi=0$), with vanishing endomorphism
    $E=0$ and bundle curvature $\Omega=0$, the first coefficients are
    \begin{align}
      a_0(x) &= 1 , \\
      a_2(x) &= \frac{1}{6} R(x) , \\
      a_4(x)
      &=
      \frac{1}{360}
      \left(
        12\nabla^2R
        +5R^2
        -2R_{\mu\nu} R^{\mu\nu}
        +2R_{\mu\nu\rho\sigma} R^{\mu\nu\rho\sigma}
      \right) .
    \end{align}

    Thus, in the decomposition
    \begin{equation}
      a_4
      =
      \frac{1}{30}\nabla^2R
      +
      \alpha_1 R^2
      +
      \alpha_2 R_{\mu\nu} R^{\mu\nu}
      +
      \alpha_3 R_{\mu\nu\rho\sigma} R^{\mu\nu\rho\sigma},
    \end{equation}
    the explicit numerical coefficients are
    \begin{equation}
      \alpha_1
      =
      \frac{5}{360},
      \quad
      \alpha_2
      =
      -\frac{2}{360},
      \quad
      \alpha_3
      =
      \frac{2}{360}.
    \end{equation}

    The universal factor $(4\pi t)^{-2}$ already appears outside the coefficient series and must not be included a
    second time in $a_4$.
    On a closed manifold the integrated $\nabla^2R$ term is a boundary term, but it is part of the local coefficient.

    In four dimensions, the quadratic curvature invariants are not all
    independent.
    The Gauss--Bonnet density
    \begin{equation}
      \mathcal{G}
      =
      R_{\mu\nu\rho\sigma} R^{\mu\nu\rho\sigma}
      -
      4 R_{\mu\nu} R^{\mu\nu}
      +
      R^2
    \end{equation}
    is topological and does not contribute to the local equations of motion.
    The chosen basis $\{R^2, R_{\mu\nu}R^{\mu\nu},
  R_{\mu\nu\rho\sigma}R^{\mu\nu\rho\sigma}\}$ is therefore equivalent,
    up to a Gauss--Bonnet term, to the commonly used basis
    $\{R^2, R_{\mu\nu}R^{\mu\nu}, \mathcal{G}\}$.

  \subsection{Pole structure in four dimensions}
    \label{subsec:pole-structure}

    Substituting the heat expansion into the definition of $\zeta_{A_g}(s)$
    shows that in $d=4$ the coefficients $a_0$ and $a_2$ generate simple poles
    at $s=2$ and $s=1$, respectively.
    Those poles describe the meromorphic zeta function away from the evaluation point $s=0$.
    They must not be confused with the divergence of the zeta-regularized effective functional at $s=0$.
    In four dimensions,
    \begin{equation}
      \zeta_{A_g}(0)
      =
      \frac{1}{(4\pi)^2}
      \int_M\dd^4x\,\sqrt{g}\,a_4(x),
      \label{eq:zeta-zero-a4}
    \end{equation}
    up to the separately treated zero-mode contribution.
    Therefore
    \begin{equation}
      \frac{\dd S_\Pi^\zeta}{\dd\log\mu}
      =
      -\zeta_{A_g}(0),
      \label{eq:zeta-running-a4}
    \end{equation}
    so zeta-scale dependence is governed by $a_4$, not $a_2$.

    The proper-time cutoff instead gives the local ultraviolet contribution
    \begin{equation}
      S_{\Pi,\mathrm{loc}}^\Lambda
      =
      -\frac{1}{2(4\pi)^2}
      \int_M\dd^4x\,\sqrt{g}\,
      \left[
        \frac{\Lambda^4}{2}a_0
        +\Lambda^2a_2
        +\log\!\left(\frac{\Lambda^2}{\mu^2}\right)a_4
      \right]
      +O(\Lambda^0).
      \label{eq:cutoff-local}
    \end{equation}
    Thus $a_0$ and $a_2$ control quartic and quadratic cutoff sensitivity, while $a_4$ controls the logarithmic
    sensitivity.

  \subsection{Origin of the Einstein tensor}
    \label{subsec:einstein-origin}

    The crucial structural point is the following.

    The coefficient $a_2(x)=R/6$ implies that the explicit cutoff contribution contains
    \begin{equation}
      \Delta S_{\Pi,\mathrm{EH}}^\Lambda
      =
      -\frac{\Lambda^2}{12(4\pi)^2}
      \int_M\dd^4x\,\sqrt{g}\,R.
      \label{eq:cutoff-EH}
    \end{equation}

    Upon metric variation, this counterterm produces the Einstein tensor,
    \begin{equation}
      \frac{\delta}{\delta g^{\mu\nu}}
      \int d^4x\,\sqrt{g}\, R
      =
      \sqrt{g}\,G_{\mu\nu} .
    \end{equation}

    The cutoff therefore supplies an Einstein-tensor contribution with a calculable power-sensitive coefficient.
    It does not arise from the zeta-scale dependence of $\zeta'_{A_g}(0)$.
    A finite local counterterm is nevertheless allowed, so the renormalized coefficient must be fixed by matching.
    For the convention $S_\Pi=+\tfrac12\log\det' A_g$, the one-scalar cutoff contribution in
    \eqref{eq:cutoff-EH} is negative.
    A positive observed Newton constant is therefore not predicted by this contribution alone.

    This power-sensitive local contribution is the part shared with Sakharov's induced-gravity
    construction~\cite{Sakharov1967}: the spectral resolution scale plays the role of an explicit ultraviolet regulator.
    The comparison does not extend to a prediction of the physical coupling, whose finite normalization and sign depend
    on the full field content and matching condition.

    Higher-derivative local contributions originate from $a_4$ and yield
    four-derivative tensors (Bach-like structures) upon variation.
    Non-local form factors arise from the finite part of the determinant
    and will be treated separately.

  \section{Metric variation and renormalized spectral stress tensor}
  \label{sec:metric-variation}

  \subsection{Definition of the renormalized spectral stress tensor}
    \label{subsec:spectral-stress}

    The renormalized projective entropy functional
    $S_\Pi^{\mathrm ren}[g;\mu]$ defines a covariant stress tensor by
    \begin{equation}
      \mathcal{T}^{\Pi,{\mathrm ren}}_{\mu\nu}
      \equiv
      -\frac{2}{\sqrt{g}}
      \frac{\delta S_\Pi^{\mathrm ren}}{\delta g^{\mu\nu}} .
    \end{equation}

    Formally, prior to renormalization, the first variation reads
    \begin{equation}
      \delta S_\Pi
      =
      \frac12 \Tr\!\left(A_g^{-1}\delta A_g\right),
    \end{equation}
    so that the response is governed by the resolvent kernel.
    After renormalization, local counterterms contribute additional
    geometric variations.

    In four dimensions, and for the minimal Laplacian $\xi=0$, write the renormalized variation as
    \begin{equation}
      \delta S_\Pi^{\mathrm ren}
      =
      \int_M\dd^4x\,\sqrt{g}\,
      \left[
        c_{\mathrm EH}^{\mathrm ren}\,G_{\mu\nu}
        +c_\Lambda^{\mathrm ren}\,g_{\mu\nu}
        +\beta^{\mathrm ren}\,B_{\mu\nu}
        +T^{\mathrm non\mbox{-}loc}_{\mu\nu}
      \right]\delta g^{\mu\nu}.
      \label{eq:full-decomposition}
    \end{equation}

    The local coefficients are renormalized matching data.
    Their cutoff-sensitive parts are constrained by~\eqref{eq:cutoff-local}, while their finite parts depend on the
    renormalization conditions and complete field content.

  \subsection{Local two-derivative sector}
    \label{subsec:two-derivative}

    The Einstein tensor contribution arises from the renormalized Einstein--Hilbert term.
    Its power-sensitive spectral contribution is the $a_2\Lambda^2$ term of~\eqref{eq:cutoff-EH}, which is proportional
    to
    \begin{equation}
      \int d^4x\,\sqrt{g}\, R .
    \end{equation}

    Its metric variation yields
    \begin{equation}
      \frac{\delta}{\delta g^{\mu\nu}}
      \int d^4x\,\sqrt{g}\, R
      =
      \sqrt{g}\, G_{\mu\nu}.
    \end{equation}

    After adding the bare and finite counterterm contributions, define
    \begin{equation}
      c_{\mathrm EH}^{\mathrm ren}
      =
      \frac{1}{16\pi G_N}.
      \label{eq:cEH-matching}
    \end{equation}
    This equation is a matching definition of the observed Newton coupling, not a prediction of its finite value.

    The cosmological term similarly receives an $a_0\Lambda^4$ contribution and an independent finite matching term.
    It cannot be discarded by an infrared approximation.

  \subsection{Local four-derivative sector}
    \label{subsec:four-derivative}

    The coefficient $a_4$ fixes the logarithmic scale dependence of quadratic-curvature counterterms of the form
    \begin{equation}
      \int d^4x\,\sqrt{g}\,
      \left(
        c_1^{\mathrm ren} R^2
      +
        c_2^{\mathrm ren} R_{\mu\nu}R^{\mu\nu}
        +
        c_3^{\mathrm ren} R_{\mu\nu\rho\sigma}R^{\mu\nu\rho\sigma}
      \right).
    \end{equation}

    The numerical coefficients in $a_4$ determine the logarithmic beta functions of these couplings, not their finite
    renormalized values.

    Their variation produces fourth-order differential tensors,
    which we collectively denote by $B_{\mu\nu}$.
    These are Bach-like structures containing four derivatives of the metric.

  \subsection{Non-local and logarithmic responses}
    \label{subsec:nonlocal-anomaly}

    The finite part of the determinant generates non-local contributions.
    In covariant form, these can be expressed through curvature-dependent
    form factors of the type
    \begin{equation}
      R \,\mathcal{F}\!\left(\frac{\Box}{\mu^2}\right) R ,
    \end{equation}
    and similar structures, as developed by
    Barvinsky and Vilkovisky~\cite{BarvinskyVilkovisky1985,BarvinskyVilkovisky1990}.
    Their metric variation yields non-local tensors
    $T^{\mathrm{non\mbox{-}loc}}_{\mu\nu}$.

    In addition, zeta regularization produces the local scale response~\cite{BirrellDavies1982,Vassilevich2003}
    \begin{equation}
      \frac{\dd S_\Pi^\zeta}{\dd\log\mu}
      =
      -\frac{1}{(4\pi)^2}
      \int_M\dd^4x\,\sqrt{g}\,a_4(x).
    \end{equation}

    Its metric variation is already included in the local four-derivative tensor $B_{\mu\nu}$ of
    \eqref{eq:full-decomposition}; it is not an additional independent tensor.
    A conformal-anomaly interpretation requires the conformal operator and its endomorphism terms; the present minimal
    operator is not conformally coupled.
    In either case, the logarithmic local response is of four-derivative order.

  \subsection{Infrared hierarchy}
    \label{subsec:ir-hierarchy}

    Consider the regime in which all dimensionless curvature invariants satisfy
    \begin{equation}
      R\,\ell_{\mathrm{sp}}^2 \ll 1,
      \qquad
      R_{\mu\nu}\,\ell_{\mathrm{sp}}^2 \ll 1.
    \end{equation}

    The derivative hierarchy is controlled by the ratio of renormalized coefficients, not by the cutoff scale alone.
    Define a higher-derivative length by
    \begin{equation}
      L_4^2
      \equiv
      \left|\frac{\beta^{\mathrm ren}}{c_{\mathrm EH}^{\mathrm ren}}\right|.
      \label{eq:L4-def}
    \end{equation}
    Since $B_{\mu\nu}$ is quadratic in curvature, its relative contribution scales as $R L_4^2$.
    If both coefficients are dominated by their spectral-cutoff contributions without cancellations, then
    $L_4=O(\ell_{\mathrm{sp}})$.
    This identification is a matching assumption, not a consequence of zeta regularization.

    Therefore, for $R L_4^2\ll1$,
    the local two-derivative Einstein term dominates over the
    four-derivative sector.
    Non-local form factors are additionally controlled by the external momentum or curvature scale relative to the
    spectral matching scale.

    This establishes a conditional infrared hierarchy once the renormalized matching coefficients are specified,
    with a non-zero Einstein coefficient.
    The hierarchy is a power-counting statement about the local sectors: it gives no general bound on the non-local
    contributions.

  \section{Spectral cutoff and Newton-coefficient matching}
  \label{sec:newton-identification}

  \subsection{Quadratic cutoff sensitivity}
    \label{subsec:running}

    For the minimal scalar Laplacian with $a_2=R/6$, the explicit proper-time cutoff contribution to the
    Einstein--Hilbert coefficient is
    \begin{equation}
      \Delta c_{\mathrm EH}^{\Lambda}
      =
      -\frac{\Lambda^2}{12(4\pi)^2}.
      \label{eq:cEH-cutoff}
    \end{equation}
    Equivalently,
    \begin{equation}
      \frac{\dd}{\dd\log\Lambda}
      \Delta c_{\mathrm EH}^{\Lambda}
      =
      -\frac{\Lambda^2}{6(4\pi)^2}.
      \label{eq:cEH-cutoff-running}
    \end{equation}
    Conflating this quadratic cutoff sensitivity with zeta-scale running drops a factor $\Lambda^2$; the two
    prescriptions must be kept separate.
    Zeta-scale dependence in four dimensions is instead governed by $a_4$ through~\eqref{eq:zeta-running-a4}.

    The renormalized coefficient has the schematic matching form
    \begin{equation}
      c_{\mathrm EH}^{\mathrm ren}
      =
      c_{\mathrm EH}^{\mathrm bare}(\Lambda)
      +\Delta c_{\mathrm EH}^{\Lambda}
      +\delta c_{\mathrm EH}^{\mathrm fin}(\Lambda,\mu).
      \label{eq:cEH-renormalized}
    \end{equation}
    Its finite value is fixed by a renormalization condition or measurement.
    Contributions from additional bosonic and fermionic operators alter both the multiplicity and the sign.

  \subsection{Spectral cutoff contribution and matching scale}
    \label{subsec:spectral-cutoff}

    Identifying the physical spectral cutoff with the inverse spectral resolution scale,
    \begin{equation}
      \Lambda=\ell_{\mathrm{sp}}^{-1},
    \end{equation}
    the magnitude of the one-scalar power-sensitive contribution is
    \begin{equation}
      \left|\Delta c_{\mathrm EH}^{\Lambda}\right|
      =
      \frac{1}{12(4\pi)^2\ell_{\mathrm{sp}}^2}.
    \end{equation}

    If the complete matched coefficient is dominated by cutoff contributions and is positive, dimensional scaling
    gives $G_N=O(\ell_{\mathrm{sp}}^2)$ with a field-content-dependent loop factor.
    Neither the sign nor the numerical coefficient follows from the single minimal scalar determinant.

  \subsection{Infrared consistency}
    \label{subsec:ir-consistency}

    The local infrared hierarchy established in Section~\ref{subsec:ir-hierarchy} requires
    \begin{equation}
      R L_4^2\ll1.
    \end{equation}
    Under cutoff-dominated matching, $L_4=O(\ell_{\mathrm{sp}})$ and this reduces parametrically to
    $R\ell_{\mathrm{sp}}^2\ll1$.
    Without that matching assumption, the derivative hierarchy must be stated directly in terms of the renormalized
    coefficient ratio $L_4$.

  \section{Ultraviolet completion: admissibility and non-uniqueness of the
  Born--Infeld structure}
  \label{sec:UV-completion}

  Under the matching conditions of Section~\ref{subsec:ir-hierarchy}, the Einstein response dominates locally for
  $RL_4^2\ll1$.
  In the ultraviolet regime
  $R\,\ell_{\mathrm{sp}}^{2} \sim 1$,
  higher-derivative corrections become comparable to the
  Einstein--Hilbert term and the expansion breaks down.
  The question is whether the full non-perturbative
  effective action is determined by the same spectral
  structure, or requires independent input.

  In the capacity analysis of~\cite{Beau2026c},
  the bounded relational capacity axiom with bound
  $c_{\mathrm{BI}}$ motivates a
  square-root action density.
  We formulate the corresponding admissibility conditions
  for the gravitational sector, inherited from the
  framework of~\cite{Beau2026m,Beau2026l}, and classify
  the completions they allow.
  The determinantal Born--Infeld form is an admissible
  spectral tensorial completion of the Einstein--Hilbert
  sector, but the conditions do not select it uniquely:
  they determine the completer only through a
  one-dimensional sector function subject to constraints at
  the origin and at the saturation boundary, and an
  explicit one-parameter family of distinct admissible
  completers realises this freedom.
  We state the additional input a uniqueness theorem would
  require.

  \subsection{Setup and admissibility conditions}
    \label{subsec:UV-setup}

    Define the mixed Ricci endomorphism
    \begin{equation}\label{eq:mixed-ricci}
      X^{\mu}{}_{\nu}
      \;=\;
      \ell_{\mathrm{sp}}^{2}\,
      g^{\mu\rho}\, R_{\rho\nu}\,.
    \end{equation}
    A local covariant completer takes the form
    \begin{equation}\label{eq:completer-form}
      \mathcal{C}(g, R)
      \;=\;
      \sqrt{-g}\;\bigl[\,\Phi(X) - 1\,\bigr]\,,
    \end{equation}
    where $\Phi$ is a smooth scalar function of the
    endomorphism~$X$.
    Lemma~\ref{lem:spectralcarrier} below shows that, under its restrictions, the possible arguments of $\Phi$ form
    a family $Y_\alpha$ of spectral carriers, of which $X$ is the member $\alpha = 0$.
    Conditions~(C3) and~(C5) below are stated for the curvature endomorphism $X$ and its eigenvalue sectors;
    Theorem~\ref{thm:uniqueness-completion} shows that, so applied, they select $\alpha = 0$.
    The density is written with $\sqrt{-g}$, the form it would take after a continuation to Lorentzian signature;
    that continuation is not performed here (Section~\ref{subsec:limitations}), and on the Riemannian section
    $\sqrt{-g}$ is read as $\sqrt{g}$.
    The eigenvalue algebra below does not depend on the signature.
    The effective gravitational action is then
    \[
      S_{\mathrm{coh}}^{\mathrm{full}}[g]
      \;=\;
      \frac{1}{16\pi G_N}\,
      \int \mathcal{C}(g, R)\,\mathrm{d}^4 x\,,
    \]
    where $G_N$ is the positive matched coupling defined in~\eqref{eq:cEH-matching}, not a value fixed by the minimal
    determinant.

    The following conditions encode the structural constraints
    that any admissible tensorial completion must satisfy.

    \begin{enumerate}
      \item[\textbf{(C1)}]
      \textit{Spectral covariance.}\;
      $\Phi(X)$ depends on $X$ only through its eigenvalues:
      $\Phi(U X U^{-1}) = \Phi(X)$ for all
      $U \in \mathrm{O}(d)$.

      \item[\textbf{(C2)}]
      \textit{Einstein infrared limit.}\;
      For weak curvature ($\|X\| \ll 1$):
      \[
        \Phi(X)
        = 1 + \tfrac{1}{2}\,\mathrm{tr}\, X
        + O(\|X\|^2)\,,
      \]
      reproducing the Einstein--Hilbert density at
      leading order.

      \item[\textbf{(C3)}]
      \textit{Algebraic Born--Infeld saturation.}\;
      $\Phi(X) \to 0$ as any eigenvalue of $(I + X)$
      approaches zero, with the remaining eigenvalues of $X$
      confined to a compact subset of $(-1,+\infty)$.
      The condition is read as requiring $\Phi$ to be defined on the whole region where $I + X$ is positive and to
      vanish on its boundary.
      The completer density vanishes at maximal admissible
      curvature rather than diverging, in accordance with
      the capacity bound
      of~\cite{Beau2026c}.

      \item[\textbf{(C4)}]
      \textit{Spectral isotropy.}\;
      $\Phi$ treats all eigenvalue directions symmetrically:
      $\Phi(X) = \varphi(\lambda_1, \ldots, \lambda_d)$
      with $\varphi$ symmetric.

      \item[\textbf{(C5)}]
      \textit{Coherence extensivity.}\;
      Define $S_{\mathrm{coh}} = \log \Phi$.
      For admissibly independent sectors $X_1$, $X_2$
      (in the sense of Axiom~A3
      of~\cite{Beau2026m}):
      $S_{\mathrm{coh}}(X_1 \oplus X_2)
      = S_{\mathrm{coh}}(X_1)
      + S_{\mathrm{coh}}(X_2)$.
    \end{enumerate}

    Conditions~(C1)--(C4) are direct consequences of the
    covariance, infrared consistency, and algebraic saturation
    properties formulated in the spectral framework of
    Sections~2--5 and in~\cite{Beau2026c}.
    Condition~(C5) is the only additional input; its
    derivation from the non-factorisability axiom~A3
    of~\cite{Beau2026m} is the subject of the following
    hypothesis.

    \begin{hypothesis}[\textbf{H-ext}: admissible coherence
    extensivity]\label{hyp:Hext}
      Let $X_1$ and $X_2$ be curvature endomorphisms
      associated with admissibly independent sectors in the
      sense of Axiom~A3 of~\cite{Beau2026m}: no
      cross-admissibility constraint couples the two sectors
      after projection locking.
      Then the coherence functional is extensive:
      \begin{equation}\label{eq:H-ext}
        S_{\mathrm{coh}}(X_1 \oplus X_2)
        =
        S_{\mathrm{coh}}(X_1) + S_{\mathrm{coh}}(X_2)\,.
      \end{equation}
      In particular, the eigenvalue sectors of the curvature
      endomorphism at a point are taken to qualify as
      admissibly independent after projection locking, so
      that extensivity applies to the rank-one spectral
      decomposition used in
      Lemma~\ref{lem:extensivity-consequence}.
    \end{hypothesis}

    \begin{remark}
      Axiom~A3 of~\cite{Beau2026m} forbids premature
      factorisation of the admissible fibre prior to projection
      locking.
      After locking, sectors that carry no mutual admissibility
      constraint contribute independently to the total
      coherence cost.
      The extensivity of $S_{\mathrm{coh}}$ is the post-locking
      counterpart of pre-locking non-factorisability:
      precisely because A3 prevents premature selection, the
      residual cost after locking decomposes additively over
      independently resolved sectors.
      This is structurally analogous to the extensivity of
      thermodynamic entropy for non-interacting subsystems.
      An analytical derivation of~[H-ext] from~A3 would
      make the classification below independent of that hypothesis; no such derivation is available, and
      [H-ext] is an open hypothesis.
      Hypothesis~[H-ext] concerns only the extensivity of the
      post-locking coherence functional.
      The classification of the spectral carriers (Lemma~\ref{lem:spectralcarrier}), which holds under the
      restriction stated in that lemma and, within its scope, leaves a trace parameter free, and the constraints on
      the sector function
      (Lemma~\ref{lem:1Dreduction}) do not use it.
    \end{remark}

  \subsection{Classification of the tensorial completions}
    \label{subsec:UV-uniqueness}

    The classification proceeds in three lemmas.
    Lemma~\ref{lem:1Dreduction} does not use [H-ext];
    Lemma~\ref{lem:spectralcarrier} does not use it either, and holds under its stated restriction on the
    carrier;
    Lemma~\ref{lem:extensivity-consequence} is conditional
    on~[H-ext].

    \begin{lemma}[Spectral decomposition from coherence
    extensivity]
      \label{lem:extensivity-consequence}
      Assume~\textup{[H-ext]}.
      Then conditions~\textup{(C1)}, \textup{(C4)},
      and~\textup{(C5)} together force
      \begin{equation}\label{eq:spectral-decomposition}
        S_{\mathrm{coh}}(X)
        = \sum_{i=1}^{d} h(\lambda_i)\,,
      \end{equation}
      for a single smooth function
      $h : (-1, +\infty) \to \mathbb{R}$, where
      $\lambda_1, \ldots, \lambda_d$ are the eigenvalues
      of~$X$.
    \end{lemma}

    \begin{proof}
      The extensivity of $S_{\mathrm{coh}}$
      implies $\Phi(X_1 \oplus X_2)
      = \Phi(X_1)\,\Phi(X_2)$ by exponentiation.
      Write $X$ in diagonal form
      $X = \mathrm{diag}(\lambda_1, \ldots, \lambda_d)$
      relative to the $g$-orthonormal eigenbasis.
      Apply the multiplicativity of $\Phi$ iteratively to
      the rank-1 decomposition
      $X = \lambda_1 P_1 \oplus \cdots \oplus \lambda_d P_d$
      (where $P_i$ are rank-1 projectors onto the
      eigenspaces):
      $\Phi(X) = \prod_{i} \Phi(\lambda_i P_i)$.
      Taking logarithms
      gives~\eqref{eq:spectral-decomposition}
      with $h(\lambda) = \log \Phi(\lambda P)$ for any unit
      projector~$P$.
      Condition~(C4) ensures that $h$ does not depend on the
      choice of eigendirection.
    \end{proof}

    \begin{lemma}[Constraints on the sector function]
      \label{lem:1Dreduction}
      Under conditions~\textup{(C2)} and~\textup{(C3)},
      the sector function $h$
      in~\eqref{eq:spectral-decomposition} satisfies
      \begin{equation}\label{eq:h-unique}
        h(0) = 0\,,
        \qquad
        h'(0) = \tfrac{1}{2}\,,
        \qquad
        \lim_{\lambda \to -1^{+}} h(\lambda) = -\infty\,,
      \end{equation}
      and, conversely, every smooth
      $h : (-1, +\infty) \to \mathbb{R}$
      satisfying~\eqref{eq:h-unique} defines a completer
      $\Phi(X) = \exp\bigl(\sum_i h(\lambda_i)\bigr)$
      obeying conditions~\textup{(C1)--(C5)}.
    \end{lemma}

    \begin{proof}
      \textit{Infrared constraint.}\;
      Condition~(C2) requires, for small $\|X\|$:
      \[
        \Phi(X)
        = \exp\!\Bigl(\sum_i h(\lambda_i)\Bigr)
        \approx 1 + \sum_i h'(0)\,\lambda_i
        = 1 + h'(0)\,\mathrm{tr}\, X\,.
      \]
      Matching with $1 + \tfrac{1}{2}\,\mathrm{tr}\, X$
      gives $h'(0) = 1/2$.
      Since $\Phi(0) = 1$ at flat space, $h(0) = 0$.

      \textit{Saturation constraint.}\;
      Condition~(C3) requires $\Phi(X) \to 0$ as
      $\lambda \to -1^+$, hence
      $h(\lambda) \to -\infty$ as $\lambda \to -1^+$.

      \textit{Sufficiency.}\;
      Conversely, let $h$ be smooth and
      satisfy~\eqref{eq:h-unique}.
      Then $\Phi(X) = \exp\bigl(\sum_i h(\lambda_i)\bigr)$
      is a symmetric function of the eigenvalues alone,
      hence obeys~(C1) and~(C4); it is multiplicative over
      direct sums, hence obeys~(C5); its expansion at small
      $\|X\|$ is
      $1 + \tfrac12\,\mathrm{tr}\,X + O(\|X\|^2)$,
      hence it obeys~(C2); it is defined wherever $I + X$ is
      positive, since $h$ is defined on $(-1, +\infty)$, and
      when one or more eigenvalues of $I + X$ approach zero
      the corresponding terms $h(\lambda_i)$ diverge to
      $-\infty$ while the remaining terms stay bounded above
      near the limiting point, so $\Phi(X) \to 0$ and~(C3)
      holds in the reading stated with it.
      Conditions~(C2) and~(C3) therefore constrain only the
      value and first derivative of~$h$ at the origin and
      its boundary behaviour at saturation; the remainder of
      the sector function is left free.
    \end{proof}

    \begin{lemma}[Spectral carrier family]
      \label{lem:spectralcarrier}
      Let the argument of~$\Phi$ be the endomorphism obtained by raising one index of a local covariant symmetric
      tensor on $(M, g)$ that is linear in the curvature and of two-derivative order, require its spectrum not to be
      a multiple of the identity in general, and normalise it so that its trace is
      $t \equiv \mathrm{tr}\,X = \ell_{\mathrm{sp}}^{2} R$, the normalisation for which~(C2) reproduces the
      Einstein--Hilbert density with the length $\ell_{\mathrm{sp}}$.
      Then the argument is a member of the family
      \begin{equation}\label{eq:carrier-family}
        Y_\alpha
        \;=\;
        \frac{\ell_{\mathrm{sp}}^{2}\, g^{-1}\bigl(\mathrm{Ric} - \alpha R\, g\bigr)}{1 - \alpha d}
        \;=\;
        \frac{t}{d}\, I + \frac{X_0}{1 - \alpha d}\,,
        \qquad
        \alpha \neq \frac{1}{d}\,,
      \end{equation}
      where $X_0 = X - (t/d)\,I$ is the traceless part of the mixed Ricci endomorphism~\eqref{eq:mixed-ricci}.
      Every member has trace~$t$, so all members induce the same linear (Einstein) term, and
      \begin{equation}\label{eq:carrier-second-order}
        \mathrm{tr}\, Y_\alpha^{2} - \mathrm{tr}\, X^{2}
        \;=\;
        \Bigl[(1 - \alpha d)^{-2} - 1\Bigr]\, \mathrm{tr}\, X_0^{2}\,,
      \end{equation}
      which vanishes for $\alpha = 0$ and for $\alpha = 2/d$, and is non-zero for every other $\alpha$ whenever
      $X_0 \neq 0$.
      Within the restriction of the lemma, the members are distinguished neither by
      conditions~\textup{(C1)}, \textup{(C2)}, \textup{(C4)} nor by the $a_2$ coefficient, and second-order data
      determine at most $|1 - \alpha d|$, leaving the pair $\{\alpha,\, 2/d - \alpha\}$.
      Conditions~\textup{(C3)} and~\textup{(C5)} are not used in the lemma; Theorem~\ref{thm:uniqueness-completion}
      shows that, applied as stated to the curvature endomorphism $X$ and its eigenvalue sectors, they select
      $\alpha = 0$.
    \end{lemma}

    \begin{proof}
      At two-derivative order, the symmetric $(0,2)$ tensors linear in the curvature are the linear combinations of
      $R_{\mu\nu}$ and $R\, g_{\mu\nu}$.
      Condition~(C1) requires the argument of~$\Phi$ to be an endomorphism $T_p M \to T_p M$, obtained by raising one
      index with the metric; it is therefore $a X + b\, t\, I$ for constants $a$ and $b$.
      For $a = 0$ the argument is $f(R)\,\delta^{\mu}{}_{\nu}$, whose spectrum is completely degenerate,
      $\lambda_1 = \cdots = \lambda_d$: the completer is then a function of $R$ alone and cannot encode the
      anisotropic curvature content.
      This case is excluded by the requirement on the spectrum, which is a condition of the lemma and not a
      consequence of~(C1)--(C4).
      The trace normalisation gives $a + b d = 1$.
      Writing $a = (1-\alpha d)^{-1}$ and $b = -\alpha (1-\alpha d)^{-1}$ parametrises every pair with $a \neq 0$
      and $a + b d = 1$, with $\alpha = -b/a$; this is~\eqref{eq:carrier-family}, and the decomposition into
      $(t/d)\,I$ and the rescaled traceless part follows from $a X + b t I = a X_0 + (a + b d)(t/d)\,I$.
      An overall rescaling of the argument that is not compensated by the trace normalisation is a redefinition of
      $\ell_{\mathrm{sp}}^{2}$.
      Since $X_0$ is traceless, $\mathrm{tr}\, Y_\alpha^{2} = t^{2}/d + (1-\alpha d)^{-2}\,\mathrm{tr}\, X_0^{2}$ and
      $\mathrm{tr}\, X^{2} = t^{2}/d + \mathrm{tr}\, X_0^{2}$, which gives~\eqref{eq:carrier-second-order}; the
      bracket vanishes exactly when $1 - \alpha d = \pm 1$.
      More generally, every second-order invariant of $Y_\alpha$ is a combination of $t^{2}$ and
      $(1-\alpha d)^{-2}\,\mathrm{tr}\,X_0^{2}$, so second-order data determine at most $(1-\alpha d)^{2}$.
      Conditions~(C1) and~(C4) concern only the dependence of~$\Phi$ on the eigenvalues of its argument, and~(C2)
      fixes the trace normalisation; none of them involves $\alpha$.
      The coefficient $a_2 = R/6$ is first-order data: it fixes the trace normalisation and not $\alpha$.
    \end{proof}

    \begin{remark}[Branches and admissibility domains of the carriers]
      \label{rem:carrier-domains}
      For $\alpha > 1/d$ the traceless part of $Y_\alpha$ changes sign, while the isotropic part $(t/d)\,I$ is the
      same for every~$\alpha$; the carrier as a whole is not opposite to the Ricci endomorphism.
      The lemma keeps both branches.
      The two members of a pair $\{\alpha,\, 2/d - \alpha\}$ have traceless parts of opposite signs and equal second
      moments; data of third order can distinguish them, through terms odd in $X_0$, but need not.

      With $\nu_i$ the eigenvalues of $X_0$, the domain of the carrier $Y_\alpha$, on which every eigenvalue of
      $I + Y_\alpha$ is positive, is
      \begin{equation}\label{eq:carrier-domain}
        D_\alpha
        \;=\;
        \Bigl\{ X :\ 1 + \frac{t}{d} + \frac{\nu_i}{1 - \alpha d} > 0 \ \text{ for all } i \Bigr\}.
      \end{equation}
      In the isotropic direction $X = \lambda I$ every factor equals $1 + \lambda$, and the boundary is
      $\lambda = -1$ for every~$\alpha$.
      Along a rank-one direction $X = \lambda P$, with $P$ a unit projector, the factor along $P$ is
      $1 + \lambda(1-\alpha)/(1-\alpha d)$, which for $\alpha \neq 1$ vanishes at
      $\lambda_{P} = -(1-\alpha d)/(1-\alpha)$ and is identically $1$ for $\alpha = 1$, and the
      $d-1$ transverse factors are $1 - \alpha\lambda/(1-\alpha d)$, which for $\alpha \neq 0$ vanish at
      $\lambda_{\perp} = (1-\alpha d)/\alpha$.
      The admissible interval of~$\lambda$ is the intersection of the two corresponding half-lines:
      $(\lambda_P, +\infty)$ for $\alpha < 0$, $(-1, +\infty)$ for $\alpha = 0$, $(\lambda_P, \lambda_\perp)$ for
      $0 < \alpha < 1/d$, $(\lambda_\perp, \lambda_P)$ for $1/d < \alpha < 1$, and $(\lambda_\perp, +\infty)$ for
      $\alpha \geq 1$.
      The boundary of $D_\alpha$ therefore moves with~$\alpha$ in every anisotropic direction, and for
      $\alpha \neq 0$ it can be reached through the transverse factors rather than through the factor along~$P$.
      Only for $\alpha = 0$ does it coincide with the saturation boundary named in condition~(C3), where an
      eigenvalue of $I + X$ vanishes.
    \end{remark}

    \begin{theorem}[Admissibility and non-uniqueness of the
    tensorial completion]
      \label{thm:uniqueness-completion}
      Let $\mathcal{C}(g, R)$ be a local covariant scalar
      density of the form~\eqref{eq:completer-form}, with the argument of $\Phi$ a carrier of
      Lemma~\ref{lem:spectralcarrier},
      satisfying conditions~\textup{(C1)--(C4)} and
      Hypothesis~\textup{[H-ext]}, with~\textup{(C3)} and~\textup{(C5)} applied, as stated, to the curvature
      endomorphism~$X$ and to its eigenvalue sectors.
      Then:
      \begin{enumerate}
        \item[\textup{(o)}]
        among the completers $\Phi = \exp\sum_i h(\mu_i)$ built on a carrier $Y_\alpha$ of
        Lemma~\ref{lem:spectralcarrier}, with $\mu_i$ its eigenvalues and $h$ obeying~\eqref{eq:h-unique}, only
        $\alpha = 0$ satisfies~\textup{(C3)} as stated: the carrier is the mixed Ricci endomorphism~$X$;

        \item[\textup{(i)}]
        the determinantal Born--Infeld density
        \begin{equation}\label{eq:unique-completer}
          \mathcal{C}_{\mathrm{BI}}(g, R)
          \;=\;
          \sqrt{-g}\;
          \Bigl[\,
          \sqrt{\det\!\bigl(
          \delta^{\mu}{}_{\nu}
          + \ell_{\mathrm{sp}}^{2}\,
          g^{\mu\rho} R_{\rho\nu}
          \bigr)}
          \;-\; 1
          \,\Bigr]
        \end{equation}
        is an admissible completion, corresponding to the
        sector function
        $h(\lambda) = \tfrac12\log(1+\lambda)$;

        \item[\textup{(ii)}]
        it is not the only one: the general admissible
        completer is
        \begin{equation}\label{eq:general-completer}
          \Phi(X)
          = \exp\!\Bigl(\sum_{i=1}^{d} h(\lambda_i)\Bigr)\,,
        \end{equation}
        with $h$ any smooth function satisfying the
        constraints~\eqref{eq:h-unique}.
        In particular, the one-parameter family
        \begin{equation}\label{eq:counterexample-family}
          \Phi_a(X)
          \;=\;
          \sqrt{\det(I + X)}\;
          \exp\!\bigl(a\,\mathrm{tr}\,X^2\bigr)\,,
          \qquad a \in \mathbb{R}\,,
        \end{equation}
        corresponding to
        $h_a(\lambda) = \tfrac12\log(1+\lambda)
        + a\lambda^2$, satisfies~\textup{(C1)--(C5)} for
        every~$a$ and differs from
        $\Phi_0 = \sqrt{\det(I+X)}$ at second order in the
        infrared expansion.
      \end{enumerate}
    \end{theorem}

    \begin{proof}
      \textit{(o)}\;
      Let $\Phi = \exp\sum_i h(\mu_i)$ on a carrier $Y_\alpha$, with $h$ smooth, hence finite, on $(-1, +\infty)$
      and divergent only as its argument tends to $-1$ (Lemma~\ref{lem:1Dreduction}).
      In an eigenbasis of $X$, $\mu_i = (\lambda_i - \alpha t)/(1 - \alpha d)$.
      Let $\lambda_1 \to -1$ while the remaining eigenvalues of $X$ range over a compact subset of $(-1, +\infty)$
      with non-empty interior, as in~(C3), and write $s = \lambda_2 + \cdots + \lambda_d$.
      The eigenvalue $\mu_1$ then tends to $(-1 + \alpha - \alpha s)/(1 - \alpha d)$, which equals $-1$ for all such
      $s$ only if $\alpha(1 - d - s) = 0$ identically, that is only if $\alpha = 0$; the remaining $\mu_j$ tend to
      values that are not $-1$ for generic $\lambda_j$.
      Condition~(C3), in the reading stated with it, requires the completer to be defined on the whole region
      where $I + X$ is positive and to vanish on its boundary.
      For $\alpha \neq 0$ the completer therefore either leaves its domain $D_\alpha$ before an eigenvalue of $I + X$
      vanishes (as it always does for $0 < \alpha < 1/d$), or tends to a finite non-zero value there, and~(C3)
      fails.
      Hence $\alpha = 0$ and the carrier is $X$.
      Since a symmetric function of the eigenvalues of $Y_\alpha$ is also a symmetric function of those of $X$, the
      carrier matters only through the form in which the conditions are imposed; with~(C3) and~(C5) imposed on $X$
      and its eigenvalue sectors, as stated, the completer is a function of the eigenvalues of $X$, and extensivity
      over those sectors, which [H-ext] asserts, is the property used in
      Lemma~\ref{lem:extensivity-consequence}.

      \textit{(i) and (ii)}\;
      By Lemma~\ref{lem:extensivity-consequence}
      (conditional on~[H-ext]),
      $\Phi(X) = \exp\!\bigl(\sum_i h(\lambda_i)\bigr)$.
      By Lemma~\ref{lem:1Dreduction}, admissibility is
      equivalent to the constraints~\eqref{eq:h-unique}
      on~$h$.
      The choice $h(\lambda) = \tfrac12\log(1+\lambda)$
      satisfies them, and gives
      \[
        \Phi(X)
        = \exp\!\Bigl(\,
        \tfrac{1}{2}\sum_i \log(1 + \lambda_i)
        \Bigr)
        = \sqrt{\,\prod_i (1 + \lambda_i)}
        = \sqrt{\det(I + X)}\,,
      \]
      which proves~(i) upon substitution
      into~\eqref{eq:completer-form}.
      For~(ii), each $h_a$ is smooth on $(-1,+\infty)$ with
      $h_a(0) = 0$, $h_a'(0) = \tfrac12$, and
      $h_a(\lambda) \to -\infty$ as $\lambda \to -1^{+}$,
      since the quadratic term stays bounded there; by the
      sufficiency part of Lemma~\ref{lem:1Dreduction},
      every $\Phi_a$ obeys~(C1)--(C5).
      For $a \neq 0$ the completers are distinct:
      $\Phi_a(X)
      = \Phi_0(X)\,
      \bigl(1 + a\,\mathrm{tr}\,X^2 + O(\|X\|^4)\bigr)$.
    \end{proof}

    \begin{remark}[Extensivity over curvature sectors and the carrier]
      Hypothesis~[H-ext], applied to the eigenvalue sectors of $X$, is satisfied by the carrier $\alpha = 0$
      (Lemma~\ref{lem:extensivity-consequence}).
      For a carrier $Y_\alpha$ with $\alpha \neq 0$ the eigenvalues $\mu_i$ depend on the total trace $t$, which
      couples the sectors: at second order, multiplicativity of $\exp\sum_i h(\mu_i)$ over the rank-one sectors of
      $X$ fails by a term proportional to $\alpha(\alpha d - 2)\,h''(0)$; where that term vanishes ($\alpha = 2/d$, or
      $h''(0) = 0$) the failure, if any, appears only at higher order.
      The selection of the carrier in Theorem~\ref{thm:uniqueness-completion} rests on~(C3); the non-uniqueness of
      the completion lies in the sector function~$h$.
    \end{remark}

    \begin{remark}[Input required for a uniqueness theorem]
      \label{rem:uniqueness-input}
      Conditions~(C1)--(C5) fix the completer only through
      the constraints~\eqref{eq:h-unique}; selecting
      $h(\lambda) = \tfrac12\log(1+\lambda)$ within the
      family~\eqref{eq:general-completer} requires input
      beyond them; the carrier is selected by~(C3) (Theorem~\ref{thm:uniqueness-completion}(o)).
      Three candidate sources of such input are open.
      First, a derived reduction of the tensorial sector
      function to the scalar saturation profile
      of~\cite{Beau2026c} would fix $h$ pointwise; at
      present no construction identifies the eigenvalue
      sectors of the curvature endomorphism with the scalar
      Born--Infeld degree of freedom.
      Second, matching the infrared expansion of $\Phi$
      beyond first order to independently computed
      higher-order heat-kernel data would eliminate
      polynomial deformations order by order; carried to
      all orders this would amount to specifying $\Phi$
      rather than deriving it, so only finitely many orders
      can legitimately serve as input; the present analysis supplies none, since the coefficient $a_4$ fixes
      logarithmic beta functions and not finite values (Section~\ref{subsec:four-derivative}).
      Third, a propagation condition on the induced field
      equations---absence of additional characteristics or
      of birefringence, in the sense in which exceptional
      propagation selects Born--Infeld electrodynamics among
      the nonlinear Lagrangians $L(S,P)$ with a Maxwell
      limit, Maxwell theory being the linear
      case~\cite{Boillat1970,Plebanski1970}---could
      single out the determinantal form; whether the
      spectral framework implies such a condition is open.
    \end{remark}

    Using $\det(g \cdot A) = \det g \cdot \det A$, the
    determinantal member takes the equivalent
    Eddington--Born--Infeld form:

    \begin{corollary}[Eddington--Born--Infeld action]
      \label{cor:EBI}
      For the determinantal
      member~\eqref{eq:unique-completer} of the admissible
      family of Theorem~\ref{thm:uniqueness-completion},
      the gravitational action takes the form
      \begin{equation}\label{eq:EBI-action}
        S_{\mathrm{coh}}^{\mathrm{full}}[g]
        \;=\;
        \frac{1}{16\pi G_N}\,
        \int_{M}
        \Bigl[\,
        \sqrt{-\det\!\bigl(
        g_{\mu\nu}
        + \ell_{\mathrm{sp}}^{2}\, R_{\mu\nu}
        \bigr)}
        \;-\;
        \sqrt{-g}
        \,\Bigr]\,
        \mathrm{d}^4 x\,.
      \end{equation}
    \end{corollary}

  \subsection{Structural interpretation}
    \label{subsec:UV-interpretation}

    \begin{remark}
      The action~\eqref{eq:EBI-action} coincides in form
      with Eddington-inspired Born--Infeld gravitational
      actions studied independently in the modified gravity
      literature~\cite{DeserGibbons1998,BanadosFerreira2010}.
      The present setting differs in two respects.
      First, the EiBI form is not postulated as a modified
      gravity model but arises as a distinguished member of
      the family of completions allowed by the admissibility
      conditions---the member, on the carrier selected by~(C3), whose sector
      function carries no deformation beyond the
      constraints~\eqref{eq:h-unique}; its selection within
      that family is an open problem
      (see the remark ``Input required for a uniqueness theorem'' following
      Theorem~\ref{thm:uniqueness-completion}).
      Second, the spectral capacity length $\ell_{\mathrm{sp}}$ specifies the physical cutoff entering the
      power-sensitive contributions.
      It does not by itself fix the finite Newton constant.
    \end{remark}

    \begin{remark}
      The framework involves two spectrally defined objects
      that must be carefully distinguished.
      At the \emph{global} level, the projective entropy
      functional
      $S_\Pi[g] = \tfrac{1}{2}\log\det' A_g$
      of Section~\ref{sec:projective-entropy-covariant}
      encodes the full eigenvalue spectrum of the effective
      Laplacian~$A_g$ and supplies the cutoff-sensitive contribution to the infrared Einstein coefficient
      (Section~\ref{sec:metric-variation}).
      At the \emph{local} level, the coherence density
      $S_{\mathrm{coh}}(X)
      = \tfrac{1}{2}\log\det(I + X)$
      encodes the eigenvalue spectrum of the mixed Ricci
      endomorphism and governs the ultraviolet saturation
      regime.
      Both share the algebraic form
      $\tfrac{1}{2}\log\det(\cdot)$ because both originate
      from the extensivity of admissible coherence:
      the former over eigenvalue modes of~$A_g$, the latter
      over eigenvalue directions of~$g^{-1}R$.
      Admissible spectral extensivity is therefore the common
      principle underlying both the infrared gravitational
      response and the ultraviolet Born--Infeld completion.
    \end{remark}

    \begin{remark}
      Lemma~\ref{lem:1Dreduction} does not use [H-ext]:
      it is an elementary analysis of
      conditions~(C2) and~(C3).
      Lemma~\ref{lem:spectralcarrier} does not use [H-ext] either:
      it is an elementary classification under its stated restriction to linear two-derivative carriers whose
      spectrum is not a multiple of the identity, and within its scope it leaves the trace parameter $\alpha$ free;
      Theorem~\ref{thm:uniqueness-completion}(o) selects $\alpha = 0$ through~(C3) applied to $X$.
      Lemma~\ref{lem:extensivity-consequence} is conditional
      on~[H-ext], whose derivation from Axiom~A3
      of~\cite{Beau2026m} is structurally motivated but
      not analytically closed.
      Theorem~\ref{thm:uniqueness-completion} is therefore
      conditional on~[H-ext] and on the restriction of Lemma~\ref{lem:spectralcarrier}; its item~(o), the selection
      of the carrier among the completers $\exp\sum_i h(\mu_i)$ built on the carriers of that lemma, with $\mu_i$
      the eigenvalues of the carrier and $h$ obeying~\eqref{eq:h-unique}, uses only~(C3) and that restriction,
      not~[H-ext].
      Within the restriction, the conditionality resides in the passage from non-factorisability
      (Axiom~A3) to extensivity of~$S_{\mathrm{coh}}$.
      The selection of the determinantal member within the
      admissible family is a separate, open problem
      (see the remark ``Input required for a uniqueness theorem'' following
      Theorem~\ref{thm:uniqueness-completion}).
    \end{remark}

  \section{Discussion and outlook}
  \label{sec:discussion}

  \subsection{Summary of the conditional structural result}
    \label{subsec:summary}

    We have constructed a covariant projective entropy functional
    \begin{equation}
      S_\Pi[g] = \frac12 \log \det\nolimits' A_g ,
    \end{equation}
    for a minimal elliptic Laplace-type operator.

    Separating a physical proper-time cutoff from zeta regularization, we showed that the renormalized metric variation
    admits the decomposition
    \begin{equation}
      \delta S_\Pi^{\mathrm ren}
      =
      \int_M\dd^4x\,\sqrt{g}\,
      \left[
        c_{\mathrm EH}^{\mathrm ren}G_{\mu\nu}
        +c_\Lambda^{\mathrm ren}g_{\mu\nu}
        +\beta^{\mathrm ren}B_{\mu\nu}
        +T^{\mathrm non\mbox{-}loc}_{\mu\nu}
      \right]\delta g^{\mu\nu}.
    \end{equation}

    The $a_2\Lambda^2$ cutoff term contributes to the Einstein coefficient, while $a_4$ governs zeta-scale dependence
    and logarithmic four-derivative terms.
    Non-local form factors and logarithmic local contributions are retained explicitly.

    For a non-zero matched Einstein coefficient and $RL_4^2\ll1$, the local two-derivative Einstein term dominates over
    the four-derivative contributions, by power counting.
    Under cutoff-dominated matching one has parametrically $L_4=O(\ell_{\mathrm{sp}})$ and
    $G_N=O(\ell_{\mathrm{sp}}^2)$, but neither the sign nor the numerical coefficient of $G_N$ follows from the single
    determinant.

    The result is therefore a structural and conditional statement about the available infrared local response.

  \subsection{Structural necessity of non-injective projection}
    \label{subsec:structural-necessity}

    The spectral analysis developed above shows that the projective entropy functional permits an Einstein response
    and supplies its power-sensitive cutoff contribution.
    It is useful to clarify under which structural conditions such an emergent
    gravitational sector can arise.

    Consider a physical description obtained through a projection
    \[
      \Pi : F \rightarrow O ,
    \]
    mapping a space of admissible configurations $F$ to a space of observable
    states $O$.
    The space $F$ is not an external substrate: it is induced by the admissibility
    structure of the observable states, with its internal organisation constrained
    by the capacity axiom of the Born\textendash Infeld companion paper~\cite{Beau2026c}.
    Such projections need not be injective: distinct admissible configurations may
    yield the same observable outcome.

    We adopt the following structural reading, which is not proved here: admissible injective projections
    possess trivial holonomy.
    On this reading, if the projection is injective, transport along closed loops in the projection
    fiber is trivial.
    The associated fiber connection therefore has vanishing curvature,
    \[
      \Omega_{\mu\nu} = 0 .
    \]

    In this case the projection cannot generate additional geometric structure in
    the effective description.
    In particular, no contribution to the effective spacetime curvature can arise
    from projection-induced fiber curvature.

    When the projection is non-injective, by contrast, the projection fiber may
    carry non-trivial curvature.
    The effective geometric response may then contain terms schematically of the
    form
    \[
      R^{\mathrm{eff}}_{\mu\nu}
      \supset
      \mathrm{tr}_{g_\Pi}\!\left(\Omega_{\mu\alpha}\Omega_\nu{}^{\alpha}\right),
    \]
    showing that curvature in the projection fiber can act as a source for
    effective spacetime curvature.

    This reading provides a structural interpretation of emergent gravitation.
    If the effective description were injective, the projection would introduce no
    additional curvature and no gravitational sector would arise from projection
    structure.
    On this reading, an emergent gravitational response requires a
    non-injective projection, in line with the no-go theorem of~\cite{Beau2026l}.


    \begin{remark}[Heisenberg non-commutativity and the supplied metric]
    The non-commutativity invoked by the Cosmochrony spectral programme has a concrete discrete setting.
    In $\mathrm{Heis}_3(\mathbb{Z}/q\mathbb{Z})$, the generating set
    $S_q = \{X^{\pm 1}, Y^{\pm 1}\}$ does not include the central generator $Z$,
    yet $[X, Y] = Z$: the central direction is reached through the bracket of the horizontal directions, which is
    the non-integrability of the horizontal distribution of $\mathrm{Heis}_3$.
    Whether this structure supplies the metric $g$ on which the present paper works is open.
    Q5a~\cite{Beau2026q5a} shows that the published admissibility form converges to the zero form, so that a
    non-zero spatial limit operator is the unestablished input [H-L].
    Under [H-L] and the open lifting hypothesis [H-lift], which Q9~\cite{Beau2026q9} does not discharge,
    Q5b~\cite{Beau2026q5b} gives the co-metric $\mathrm{diag}(-A_\tau, A_H, A_H, 0)$ of rank three, with an empty
    central slot and a stipulated temporal part.
    No value of $A_\tau$, $A_H$ or of a central coefficient is established: Q8~\cite{BeauQ8} shows that the
    commutator supplies no central coefficient and that, for a supplied positive rank-three spatial form invariant
    under the full $\mathfrak{su}(2)$ action, $A_H = A_Z = 2\lambda$ with $\lambda$ free, Q10~\cite{Beau2026q10}
    derives no value, and Q11~\cite{Beau2026q11} shows that spatial $\mathrm{SU}(2)$ invariance leaves $A_\tau$
    free.
    A four-dimensional co-metric $\mathrm{diag}(-A_\tau, 2\lambda, 2\lambda, 2\lambda)$ follows only for a supplied
    rank-four extension satisfying the hypotheses of Q11, Corollary~6.1, and its signature is a hypothesis of that
    construction; the convergence of the admissibility weights, Hypothesis~[H-w], remains open~\cite{Beau2026w1}.
    The availability of the spectral Einstein sector established here therefore does not depend on this
    construction, and its reading as the response of an emergent metric is conditional on [H-L], [H-lift], a
    supplied rank-four extension and a continuation from Riemannian signature that is not performed.
    Its positive finite normalization remains a separate matching problem.
    \end{remark}

    Within the spectral framework considered in this work, the geometric response
    arises from the renormalized variation of the functional
    \[
      S_\Pi[g] = \frac12 \log \det{}' A_g .
    \]
    The heat-kernel analysis identifies the Einstein tensor among the local metric responses and computes its
    quadratic cutoff sensitivity.
    The spectral construction therefore provides the ingredients for an effective dynamics compatible with a
    non-injective projective structure, conditional on renormalized matching.

    Taken together, these results suggest the following structural picture.
    Non-injective projection allows the emergence of non-trivial fiber curvature,
    which in turn enables an effective geometric response.
    The spectral entropy functional provides a minimal mechanism for an Einstein-sector contribution.
    If the matched Einstein coefficient is non-zero and $RL_4^2\ll1$, with $L_4$ defined in~\eqref{eq:L4-def}, the
    Einstein term dominates the local four-derivative contributions in the infrared, by power counting; this gives no
    general bound on the non-local terms.
    This mechanism is consistent with the relational spectral framework
    of~\cite{Beau2026a}, in which spacetime geometry emerges operationally from
    relational configurations under non-injective projection.
    Gauge dynamics is a separate question: Q12~\cite{Beau2026q12} derives Yang\textendash Mills dynamics from the
    same functional for a supplied compact gauge group, which is not identified from the projection.

  \subsection{Relation to observer-indexed space-time}
    \label{subsec:observer-indexed}

    \emph{This subsection is interpretation, not a result.}
    The present paper works on a supplied metric; on the conditional reading in which that metric is emergent, the
    reading has a direct foundational counterpart.
    Working from extended Wigner-friend no-go theorems, Pienaar shows that rejecting the absoluteness of
    observed events forbids embedding all events in a single common background space-time, and that
    space-time notions must instead be indexed to each observer through a ``relative locality''
    condition~\cite{Pienaar2026}.
    On that reading the metric $g_{\mu\nu}$ would not be fundamental but an effective structure valid only within a
    projectable regime; identifying it with a local projective ordering rate is a further interpretation, not
    established here.
    A single global manifold shared by all observers would then not be expected: distinct effective
    descriptions define distinct projectable patches, and their agreement is a regime-dependent property
    rather than a prior structural assumption.
    What Pienaar reaches as a dilemma --- reductio or a new route to quantum gravity --- appears here as
    the second horn: the observer-indexed fragmentation of space-time would be the expected signature of an
    emergent, projection-defined geometry, with the observer-independent content carried by the
    pre-geometric admissibility structure rather than by a common metric manifold.

  \subsection{Relation to induced gravity}
    \label{subsec:induced-gravity}

    The cutoff calculation shares the heat-kernel mechanism of Sakharov's induced gravity:
    integrating spectral modes produces a quadratically cutoff-sensitive local Einstein--Hilbert contribution.
    It does not show that the observed gravitational coupling emerges from the minimal determinant.
    The cutoff fixes the magnitude of this one-operator contribution, while the sign and finite renormalized coefficient
    remain matching data controlled by the complete operator content and renormalization condition.

    No fundamental gravitational action is assumed at the starting point.
    Within the explicit cutoff scheme, the minimal spectral determinant produces a calculable local
    Einstein-tensor contribution to the effective metric response, without fixing the physical Newton coupling.

  \subsection{Limitations}
    \label{subsec:limitations}

    Several limitations must be emphasized.

    First, the derivation is performed in a Euclidean Riemannian framework, on a supplied metric.
    Although the variational structure of the spectral functional is well defined
    in this setting, a complete treatment of causal propagation requires a
    Lorentzian formulation of the theory, and reading the metric as emergent requires the inputs listed in
    Section~\ref{sec:introduction}.

    Second, non-local contributions generated by the spectral expansion are not
    eliminated.
    While they are suppressed in the infrared regime considered here,
    they may become relevant at higher curvature or shorter scales.

    Third, the finite value and sign of the Newton constant are matching data.
    The minimal scalar cutoff contribution is negative with the determinant convention used here, and additional
    operator content or a finite renormalization condition is required to obtain the observed positive coefficient.
    The classification of Section~\ref{sec:UV-completion} constrains the
    functional form of this completion but fixes neither the member of the
    admissible family nor its overall normalisation independently of the
    matter content.

    Fourth, Hypothesis~[H-ext]
    (Section~\ref{subsec:UV-setup}) is not
    derived analytically from the Foundation axioms.
    The spectral classification of the tensorial completions is conditional
    on this hypothesis, and even under it the admissibility conditions leave
    the Born--Infeld member non-uniquely selected within the
    family~\eqref{eq:general-completer}, through its sector function; the carrier is selected by condition~(C3)
    applied to the curvature endomorphism.

    Finally, the present work focuses on the emergence of the effective
    gravitational field equations.
    Questions of dynamical stability, cosmological evolution, and coupling to
    additional matter sectors lie beyond the scope of the present analysis.

  \subsection{Remark on geometric flows and spectral monotonicity}
    \label{subsec:ricci-flow}

    The functional $S_\Pi$ depends on the full eigenvalue spectrum of the
    elliptic operator $A_g$.
    Under geometric evolution equations such as the Ricci flow,
    spectral quantities associated with the Laplacian exhibit
    monotonic behavior.

    In particular, monotonicity of the first eigenvalue
    $\lambda_1(-\nabla_g^2)$ of the scalar Laplacian under Ricci flow
    has been established by Ma~\cite{Ma2006}.
    Related results for modified Laplace-type operators involving
    curvature terms were obtained by Cao~\cite{Cao2007}.
    These developments are rooted in the entropy framework introduced
    by Perelman~\cite{Perelman2002}.

    Although $S_\Pi$ is not identical to Perelman's entropy functional,
    its spectral definition suggests a structural relation between
    geometric flows and spectral entropy variation.
    Clarifying this relation would require a dedicated analysis
    and lies beyond the scope of the present work.

  \subsection{Outlook}
    \label{subsec:outlook}

    The analysis presented here shows that a minimal spectral operator supplies an $a_2\Lambda^2$ Einstein--Hilbert
    contribution under an explicit spectral cutoff.
    If the renormalized matching coefficients satisfy the stated sign and hierarchy conditions, the Einstein response is
    positive and dominates the local four-derivative contributions in the infrared, by power counting; the non-local
    terms are not bounded in general.

    Several directions follow naturally from this result.

    \begin{itemize}

      \item Extension of the spectral framework to Lorentzian signature,
      allowing a fully causal formulation of the effective dynamics and
      the study of propagating gravitational modes.

      \item Coupling of the spectral sector to additional operator
      structures describing matter degrees of freedom, leading to a
      unified description of geometric and material contributions to the
      effective stress tensor.

      \item Investigation of regimes in which higher-order curvature
      contributions become relevant.
      Section~\ref{sec:UV-completion} identifies the Born--Infeld density
      as an admissible non-perturbative completion within an explicitly
      parametrised family; the resulting strong-curvature dynamics, its
      dependence on the family parameter, and its predictions for
      gravitational collapse, bouncing cosmologies, and singularity
      resolution remain to be explored.

      \item Exploration of the relation between projection structure and
      emergent geometry, clarifying how non-injective effective
      descriptions may generate curvature responses in the infrared
      regime.

      \item Construction of the chain from the admissibility data to the supplied metric.
      Its open steps are a non-zero spatial limit operator ([H-L]; Q5a~\cite{Beau2026q5a} gives a zero-form
      limit), the lifting hypothesis [H-lift] (Q5b~\cite{Beau2026q5b}; not discharged by Q9~\cite{Beau2026q9}),
      a rank-four extension of the rank-three co-metric of Q5b, every metric coefficient value
      (Q8~\cite{BeauQ8}, Q10~\cite{Beau2026q10}, Q11~\cite{Beau2026q11}), and the continuation relating the
      Riemannian metric used here to a Lorentzian one.

    \end{itemize}

    \textbf{Interpretive outlook.}
    The result keeps open the possibility that geometric gravity is a low-energy manifestation of spectral operator
    structure, while isolating the missing step: an independent matching principle must select the positive finite
    Einstein coefficient rather than merely permit it.

    \clearpage
    \bibliography{cosmochrony-bibliography,external-refs}

\end{document}
